\section{利用滚动得到玫瑰线}
\markboth{41. 利用滚动得到玫瑰线}{41. 利用滚动得到玫瑰线}

玫瑰线是变幅内摆线和变幅外摆线的特殊情形。要证明这个事实并不困难。先看变幅内摆线，在第 35 节的推导过程中，我们取 $h=R-mR=R-r$ 的特殊情形。如图 41.1，设定圆半径为 $R$，动圆半径为 $r$，$\frac{r}{R}=m$，$O_1M=OO_1=R-r=h$，$\angle MO_1N=t$，$\angle xON=mt$，若设 $O$ 点为极点，$Ox$ 轴为极轴，$M$ 点的极坐标为 $(\rho,\theta)$，则

\noindent
\begin{minipage}[c]{0.54\textwidth}
\begin{align*}
  \rho = OM &= 2OO_1\cos\angle O_1OM \\
            &= 2(R-r)\cos\frac{t}{2} = 2h\cos\frac{t}{2}.
\end{align*}
\end{minipage}%
\hfill
\begin{minipage}[c]{0.43\textwidth}
  \centering
  \begin{tikzpicture}[scale=0.82, >=Stealth, line cap=round, line join=round]
    \def\R{1.50}
    \def\r{0.52}
    \def\mt{58}
    \def\tang{82}
    \pgfmathsetmacro{\h}{\R-\r}
    \coordinate (O) at (0,0);
    \coordinate (O1) at ({\h*cos(\mt)}, {\h*sin(\mt)});
    \coordinate (N) at ({\R*cos(\mt)}, {\R*sin(\mt)});
    \coordinate (M) at ($ (O1) + ({\mt-\tang}:\h) $);

    % 坐标轴
    \draw[->] (-1.75, 0) -- (1.85, 0) node[below, font=\scriptsize] {$x$};
    \draw[->] (0, -1.75) -- (0, 1.85) node[left, font=\scriptsize] {$y$};

    % 定圆与动圆（r = 0.52 < h = R - r = 0.98，非奇异态，动圆既不过 O 也不过 M）
    \draw[thick] (O) circle (\R);
    \draw[thick, curveblue] (O1) circle (\r);

    % 连线（OO1 = O1M = h 构成严格等腰三角形）
    \draw (O) -- (N);
    \draw (O1) -- (M);
    \draw[thick] (O) -- (M);

    % 角度弧（mt、t、theta 彻底分离无压线）
    \draw (0.28, 0) arc (0:\mt:0.28);
    \node[font=\scriptsize, fill=white, inner sep=0.4pt] at (38:0.46) {$mt$};

    \draw ($(O1)+(\mt:0.22)$) arc (\mt:{\mt-\tang}:0.22);
    \node[font=\scriptsize] at ($(O1)+({\mt-\tang/2}:0.36)$) {$t$};

    \pgfmathsetmacro{\th}{\mt-\tang/2}
    \draw (0.72, 0) arc (0:\th:0.72);
    \node[font=\scriptsize] at ({0.5*\th}:0.94) {$\theta$};

    % 标点与标签（O1 置于动圆内部左侧纯空白区，远离 y 轴与蓝圆）
    \fill (O) circle (1.2pt) node[below left=-1pt, font=\scriptsize] {$O$};
    \fill (O1) circle (1.3pt) node[left=1.5pt, font=\scriptsize] {$O_1$};
    \fill (N) circle (1.2pt) node[above right=-1pt, font=\scriptsize] {$N$};
    \fill (M) circle (1.3pt) node[below=1pt, font=\scriptsize] {$M$};
  \end{tikzpicture}
  \captionof*{figure}{图\ 41.1}
\end{minipage}

因为 $\theta=\angle xOM=mt-\frac{t}{2}$，所以 $\frac{t}{2}=\frac{\theta}{2m-1}$。记 $b=\frac{1}{2m-1}$，得 $\rho=2h\cos b\theta$。再旋转极轴，使新坐标系中的极角 $\theta_1$ 由下式决定：
\[
  \theta_1=\theta+\frac{\pi}{2b},
\]
就使曲线方程化成
\[
  \rho=2h\sin b\theta_1,
\]
可见曲线是玫瑰线。由于这时 $0<m<1$，所以从 $b=\frac{1}{2m-1}$ 得到对应的 $|b|>1$。

\noindent
\begin{minipage}[c]{0.43\textwidth}
  \centering
  \begin{tikzpicture}[scale=0.78, >=Stealth, line cap=round, line join=round]
    \def\R{0.75}
    \def\r{0.45}
    \pgfmathsetmacro{\h}{\R+\r}
    \def\mt{66}
    \def\tang{88}
    \coordinate (O) at (0,0);
    \coordinate (O1) at ({\h*cos(\mt)}, {\h*sin(\mt)});
    \coordinate (M) at ($ (O1) + ({\mt+180+\tang}:\h) $);

    % 坐标轴
    \draw[->] (-1.00, 0) -- (2.05, 0) node[below, font=\scriptsize] {$x$};
    \draw[->] (0, -1.00) -- (0, 1.80) node[left, font=\scriptsize] {$y$};

    % 圆 O 与圆 O1 (外切)
    \draw[thick] (O) circle (\R);
    \draw[thick, curveblue] (O1) circle (\r);

    % 连线
    \draw (O) -- (O1) -- (M);
    \draw[thick] (O) -- (M);

    % 角度弧（与正文 theta = mt + t/2 - pi/2 严格一致）
    \draw (0.32, 0) arc (0:\mt:0.32);
    \node[font=\scriptsize, fill=white, inner sep=0.3pt] at (43:0.48) {$mt$};

    \pgfmathsetmacro{\th}{\mt+\tang/2-90}
    \draw (0.95, 0) arc (0:\th:0.95);
    \node[font=\scriptsize] at ({0.5*\th}:1.15) {$\theta$};

    \draw ($(O1)+({\mt+180}:0.22)$) arc ({\mt+180}:{\mt+180+\tang}:0.22);
    \node[font=\scriptsize] at ($(O1)+({\mt+180+\tang/2}:0.34)$) {$t$};

    % 标点与标签
    \fill (O) circle (1.2pt) node[below left=-1pt, font=\scriptsize] {$O$};
    \fill (O1) circle (1.3pt) node[above=1pt, font=\scriptsize] {$O_1$};
    \fill (M) circle (1.3pt) node[right=1pt, font=\scriptsize] {$M$};
  \end{tikzpicture}
  \captionof*{figure}{图\ 41.2}
\end{minipage}%
\hfill
\begin{minipage}[c]{0.54\textwidth}
再看变幅外摆线，在第 34 节的推导过程中，我们取 $h=R+mR=R+r$ 的特殊情形。如图 41.2，$O_1M=OO_1=R+r=h$，$\angle MO_1O=t$，$\frac{r}{R}=m$，$\angle xOO_1=mt$，若设 $O$ 为极点，$Ox$ 为极轴，$M$ 点的极坐标为 $(\rho,\theta)$，则
\end{minipage}
\[
  \rho = OM = 2OO_1\sin\frac{t}{2} = 2h\sin\frac{t}{2}.
\]

因为 $mt-\theta=\frac{\pi}{2}-\frac{t}{2}$，所以
\[
  \theta=mt+\frac{t}{2}-\frac{\pi}{2},\qquad \frac{t}{2}=\frac{\pi+2\theta}{2(2m+1)},
\]
记 $b=\frac{1}{2m+1}$，得 $\rho=2h\sin\left(\frac{b\pi}{2}+b\theta\right)$，再旋转极轴，使新坐标系中的极角 $\theta_1=\frac{\pi}{2}+\theta$，就使曲线方程化成
\[
  \rho=2h\sin b\theta_1,
\]
所以这时也得到玫瑰线。由于现在 $m>0$，所以从 $b=\frac{1}{2m+1}$ 得到 $0<b<1$。

因而最后结论是：玫瑰线 $\rho=a\sin b\theta$，当 $0<b<1$ 时是变幅外摆线，$b>1$ 时是变幅内摆线。当 $b=1$ 时曲线化为圆周 $\rho=a\sin\theta$。

例如，最常见的三叶玫瑰线（$b=3$）和四叶玫瑰线（$b=2$）都是变幅内摆线。
